\section{Gestión del Alcance}
El alcance dice qué se hará y qué no. Se escribió desde el inicio para que, si
algo se atrasa, el equipo sepa qué puede ceder (Tabla~\ref{tab:alcance}).

\begin{table}[tbp]
	\caption{Alcance del Proyecto}
	\label{tab:alcance}
	\small\setlength{\tabcolsep}{3pt}
	\begin{tabular}{P{7.4cm}P{7.4cm}}
		\toprule
		Incluye & Queda fuera \\
		\midrule
		Un acceso peatonal piloto, en el turno vespertino. & Reconocimiento facial y tarjetas NFC. \\
		QR que cambia en el celular del estudiante, con su generación y validación. & Pase de lista dentro de los salones (asistencia por materia). \\
		Lector con cámara y torniquete a escala. & Registro de salidas y control de visitantes. \\
		Servidor que valida y registra cada intento en una base de datos. & Instalación en todos los accesos y torniquete industrial. \\
		Panel web de consulta para el personal. & Retirar la revisión manual durante el piloto. \\
		Medición del acceso antes y durante el piloto. & \\
		\bottomrule
	\end{tabular}
	\tablenote{Fuente: Ficha Inicial del Proyecto y presentación del avance.}
\end{table}

Durante el piloto la revisión manual sigue funcionando. Quien no pueda
mostrar su código entra como hasta ahora.

\subsection{Entregables}
Al cierre, el equipo entregará los productos de la Tabla~\ref{tab:entregables}.
Ninguno se da por terminado sin pasar su criterio de aceptación, y cada uno
tiene un responsable de la tabla de roles (Tabla~\ref{tab:roles}).

\begin{table}[tbp]
	\caption{Entregables del Proyecto}
	\label{tab:entregables}
	\small\setlength{\tabcolsep}{3pt}
	\begin{tabular}{P{1cm}P{7cm}}
		\toprule
		N.º & Entregable \\
		\midrule
		1 & Informe de diagnóstico \\
		2 & Prototipo del torniquete con lector \\
		3 & Servidor y base de datos \\
		4 & Página del estudiante con su QR \\
		5 & Panel web de consulta \\
		6 & Código fuente en un repositorio \\
		7 & Manual de uso y aviso de privacidad \\
		8 & Informe de resultados: antes frente a durante el piloto \\
		\bottomrule
	\end{tabular}
	\tablenote{Fuente: presentación del avance. La ficha inicial menciona los mismos productos: prototipo, informe de diagnóstico, código fuente, manual de uso, propuesta de aviso de privacidad e informe final.}
\end{table}
